\documentclass{article}
\usepackage{xcolor}
\usepackage{stepdiff}
\begin{document}

\[
\begin{stepdiff}[theme=soft]
\step{x = 1}
\step[reason={soft theme}]{x = 2}
\end{stepdiff}
\]

\[
\begin{stepdiff}[theme=minimal]
\step{x = 1}
\step[reason={minimal theme}]{x = 2}
\end{stepdiff}
\]

\[
\begin{stepdiff}[highlight=background]
\step{a=b}
\step[reason={background}]{a=c}
\end{stepdiff}
\]

\[
\begin{stepdiff}[highlight=underline]
\step{a=b}
\step[reason={underline}]{a=c}
\end{stepdiff}
\]

\[
\begin{stepdiff}[highlight=color]
\step{a=b}
\step[reason={color}]{a=c}
\end{stepdiff}
\]

\[
\begin{stepdiff}[highlight=none]
\step{a=b}
\step[reason={none}]{a=c}
\end{stepdiff}
\]

\[
\begin{stepdiff}[layout=compact]
\step{S_n = 1 + \cdots + n}
\step[reason={compact}, diff=all]{S_n = \frac{n(n+1)}{2}}
\end{stepdiff}
\]

\[
\begin{stepdiff}[layout=lecture]
\step{a_n \le b_n}
\step[reason={lecture layout}]{\lim a_n \le \lim b_n}
\end{stepdiff}
\]

\[
\begin{stepdiff}[layout=wide]
\step{P \Rightarrow Q}
\step[reason={wide layout}]{P \wedge R \Rightarrow Q}
\end{stepdiff}
\]

\[
\begin{stepdiff}[theme=minimal, layout=wide, highlight=color]
\step{x \approx y}
\step[reason={relation still aligns}, diff=false]{x \equiv y}
\step[reason={emphasis}, diff=all]{x = y}
\end{stepdiff}
\]

\end{document}
